\documentclass{article}
\usepackage{amsmath}
\usepackage{graphicx}
\usepackage{float}

\usepackage[utf8]{inputenc}     % pdfLaTeX
\usepackage[T1]{fontenc}
\usepackage[magyar]{babel}

\usepackage{xfrac}
\usepackage{nicefrac}
\usepackage{multirow}
\usepackage[top=2cm, bottom=2cm, left=3cm, right=3cm]{geometry}

\title{Árpád ház}
\author{Haydu Lőrinc, Lengyel Zoltán Péter, Rovenszky Robin, Simon László}
\date{Legutóbb frissítve: 2026. Június 08.}

\newtheorem{remark}{Megjegyzés}

\begin{document}

\maketitle
\tableofcontents

\section{Trónharc, trónviszály}
\begin{itemize}
    \item I. István halála előtt Vazult megvakíttatta, hogy ne legyen alkalmas a trónra
    \item Vazul fiait elűzték
    \item Orseolo Péter $1038-1041$. (István nővérét veszi feleségül, velencei dózse fia) 
    \begin{itemize}
        \item Németellenes
        \item Pogányüldöző
        \item Pécsi székesegyház
        \item István halála után Gizellát kisemmizi, megfosztja vagyonától és száműzi vidékre
        \item Idegeneket hív be királyi udvartartásra
        \item Aba Sámuel elűzi a trónról (István tesvérével házasodott)
    \end{itemize}
        
    \item Aba Sámuel $1041-1044$
    \begin{itemize}
        \item Német-Római császár és Orseolo Péter
        \item 1042. Osztrák Őrgrófságra betört $\rightarrow$ III: Henrik császár folyamatos támadásai
        \item Szent István politikájával szakít (nem üldözi a pogányokat)
        \item Balkáni bogumilok beengedése
        \item Parasztokkal tölti a szabadidejét
        \item Összeesküvések és III. Henrik császár, ménfői csata $\rightarrow$ elűzik, Orseolo visszakerül a trónra
    \end{itemize}
    
    \item Orseolo Péter (másodjára) $1044-1046$
    \begin{itemize}
        \item németek hűbérese
        \item bajor törvény 
        \item adók újbóli kivetése
        \item Politikája nem változott, nem tanult $\rightarrow$ három összeesküvés (Vata-féle pogánylázadás $\rightarrow$ kereszténypapok mészárlása 
        \begin{itemize}
            \item Gellért püspök (csanádi püspök) és Kelen-hegy
        \end{itemize}
        
        \item \textbf{Mártírhalál/Vértanú}: Meghal valamilyen fontos ügyért. A mártír feláldozza magát, a vértanú nem feltétlenül, csak fontos ügyért hal meg. Keresztény kifejezések. 
        
    \end{itemize}
    \item I. András $1046-1060$.
    \begin{itemize}
        \item \textbf{Dukátus} Bélával: Vagyis Béla megkapta az ország $\nicefrac{1\,}{\,3}\,$-át.
        \item Felesége Anasztázia (orosz támogatás)
        \item Eredetileg öccsét akarta örökösnek (Béla), de később fia született (Salamon).
        \item Korona vagy kard 
        \begin{itemize}
            \item Kard - lovagság
            \item Korona - királyság, DE! lemészárolta volna öccsét
        \end{itemize}
        \item Bélával való csatában meghal (Tihanyi apátságba temetik el) 
        \begin{itemize}
            \item Tihanyi apátság alapítólevele: 1055, legrégebbi magyar szórványemlék
        \end{itemize}
    \end{itemize}
    
    \item I. Béla $1060-1063$

    \item 11. évig nincs folytonos uralkodó
    
    \item I. Géza $1074-1077$
    \begin{itemize}
        \item Salamon unokatestvére
        \item VII. Gergely pápa nem támogatja $\rightarrow$ bizánci császár elismerése kell (Korona alsó része) 
        \item Vásártartás napját áthelyezi vasárnapról szombatra 
        \item Salamont nem ismeri el örököseként (csak Lászlót)
        \item Salamonnal való tárgyalás közben meghal
        \item Jövőhét hétfő
    \end{itemize}
\end{itemize}

\section{I. Szent László}
\begin{itemize}
    \item I. Szent László $1077-1095$
    \begin{itemize}
        \item Konszolidáció
        \item Béla kisebbik fia
        \item $1083$ - szentté avatta a következőket: István, Imre herceg, Gellért püspök
        \item $1091.$ Zágrábot visszafoglalta, egyházmegyét hoz létre
        \item Védi az országot a kunokkal, besenyőkkel szemben
        \item Törvényei: 
        \begin{itemize}
            \item Védte a magántulajdont
            \item Védte a kereszténységet
            \item Pogányokat üldözte
            \item Boszorkányokat üldözte 
            \begin{itemize}
                \item Tűzpróba (forró vasat döftek, ha a seb begennyesedett vagy begyógyult boszorkány)
                \item Vízpróba (Lábához nehezéket kötöttek és vízbe dobták, ha kiszabadul és feljön a víz felszínére boszorkány, ha nem akkor meg már úgyis mindegy)
            \end{itemize}
            \item Istenítélet (tűzpróba vízpróba, "Isten dönt a sorsáról", sokszor nem élték túl, ha túlélték pedig boszorkánynak számítottak)
            \item Ha lop $\rightarrow$ akasszák fel/vakítsák meg/istenítélet \begin{itemize}
                \item Ezek hatása nagyrészt ugyanaz volt, mert se a vakítást se az istenítéletet nem élték túl
            \end{itemize}
            \item verés = tanítás
            \item Akik kutak mellett áldoznak, kövekhez ajándékot visznek (pogány szokás szerint) ökörrel fizetnek
        \end{itemize}
        \item fia nem született, lányai diplomáciai érdekből ki lesznek házasítva (Piroska $\rightarrow$ bizánci szent)
    \end{itemize}
    \item XI. század $\rightarrow$ átalakuló társadalom

    \begin{tabular}{l l l}
    szabadok 
      & $\rightarrow$ ispán 
      & \multirow{3}{*}{$\left.\begin{array}{l}
          \text{várispán (comes)} \\
          \text{}
        \end{array}\right\} \text{trónharcok $\rightarrow$ nagybirtokos}$} \\
        \vspace{3pt}
     & $\rightarrow$ várjobbágy / vitéz & \\
     & $\rightarrow$ köznép & $\rightarrow$ kóborló (peremterület) \\
     & & $\rightarrow$ lesüllyedt szabad 
    \end{tabular}
    
    \item Sírja Nagyvárad
    \item Hermája - Győr
    \item Kar és állkapocs ereklye
    \item Szobra - hősök tere
    \item Freskója - Szent László templom (Béke tér)
    
\end{itemize}

\section{Könyves Kálmán}
% todo: make sure (with Gaár) that this section is OK once we are finished with it
\begin{itemize}
    \item Coloman the Learned
    \item $1095-1116$ Uralkodott
    \item Enyhít László szigorán
    \item Művelt király
    \item Könyves - gúnynév, ragadványnév 
    \item Belpolitika
    \begin{itemize}
        \item Birtokállomány helyreállítása
        \begin{itemize}
            \item háramlási jog
        \end{itemize}
        \item Kóborlók letelepítése
        \item Katonaállítási kötelezettség
        \begin{itemize}
            \item nagybirtokosok: páncélos vitéz felfegyverzése
        \end{itemize}
    \end{itemize}
    \item Külpolitika
    \begin{itemize}
        \item Sikertelen külpolitika, nem tud külterületet foglalni
        \item Dalmácia területe: sikeres
        \item Halics: sikertelen
        \item Horvát bánsággal perszonálunió
        \item Autonóm maradt Horvátország
        \item Cölibátus visszavezetés: papi nőtlenség
    \end{itemize}
\end{itemize}

\section{Kálmán again}
\begin{itemize}
    \item Magyar krónikák szerint csúnya, púpos
    \item Külföldi krónikák szerint művelt, okos
\end{itemize}
Belpolitikája:
\begin{itemize}
    \item Kóborlókat letelepíti
    \item Háromlási jog bevezetése (ha nincs ki a földet örökölje visszaszáll a királyra, királyi regálé része)
    \item királyi regálé: Királyi felségjogon szerzett jövedelem
    \item kereskedelmi adót növeli
    \item királyi regálé nő $\rightarrow$ nő a király vagyona, földbirtokainak száma
    \item kisebb értékűek a pénzek
\end{itemize}
Kálmán törvényei: 
\begin{itemize}
    \item enyhébbek mint László törvényei
    \item egységesebb
    \item nincsenek boszorkány $\rightarrow$ nincs boszorkányüldözés \begin{itemize}
        \item DE: van kuruzsló, jövendőmondó, javasasszony $\rightarrow$ TILTOTT
    \end{itemize}
\end{itemize}
Külpolitika:
\begin{itemize}
    \item Horvát királyság területeit elfoglalja
    \item 1102-1918 perszonálunió
    \item dalmát városok: Zára, Trau, Split
    \item De: területre igényt tart Velencei köztársaság, Bizánc
\end{itemize}

Egyház: 
\begin{itemize}
    \item cölibátus
\end{itemize}

\subsection{Szent korona}
\begin{itemize}
    \item Koronázási ékszer
    \item Hartvik legenda: II. Szilveszter pápától származik, I. Istvánt koronázták meg vele
    \item Két részből áll 
    \begin{itemize}
        \item Alsó abroncs - görög korona 
        \begin{itemize}
            \item Bizánc, női fejék
            \item Központi részén Krisztus, Bizánci császárok, Géza fejedelem, görög szentek
        \end{itemize}
        
        \item Felső keresztpántos rész - latin korona 
        \begin{itemize}
            \item Apostolok díszítik
        \end{itemize}
        
        \item Egyesítés feltehetően III. Béla idején (Istvánt csak a felsővel koronázták)
        \item Tetején kereszt, eredetileg egyenes de elferdült
    \end{itemize}
    
    \item 1440. Szilágyi Erzsébet szolgálójával raboltatja el $\rightarrow$ III. Frigyeshez kerül $\rightarrow$ Mátyásnak $80000$ arany forintba kerül visszavásárolniaű
    
    \item Volt Bécsben, Prágában, Pozsonyban, Budán, Orsovai mocsárban, II. vh alatt benzineshordóban Ausztriában, 1953 USA-ban kerül, 1978-ban kerül vissza hozzánk 
    
    \item 2000-ig a Nemzeti Múzeumban, azóta Parlamentben

    \item XVI. század: Szent Korona tan (A Szent Korona a királyi hatalmat testesíti meg)

    \item Utolsó koronázás 1916. IV. Károly, Ferenc József halála után

\end{itemize}



\subsection{Koronázási ékszerek}
\begin{itemize}
    \item Koronázási palást
    \item Országalma
    \item Koronázási kard
    \item Jogar
    \item Szent Korona (Hatalom megszerzése $\rightarrow$ hatalom)
\end{itemize}

\section{Fogalmak}

\begin{description}
\item[Mainzi szertartásrend:]
Koronázási ékszerekkel koronázzuk a királyt.

\item[Invesztitúra:]
A pápa megválasztásának joga.

\item[Dukátus:]
Hercegség; a király által hercegnek adott, a herceg által irányított terület.

\item[Őstörténet:]
Egy nép történetének korai szakasza.

\item[Életfa:]
Pogány szimbólum/jelkép, az életet és a túlvilágot jelképezi.

\item[Királytükör:]
A király megfogalmazza benne uralkodási tanácsait.

\item[Várjobbágy:]
Katona.

\item[Primogenitúra:]
Elsőszülött fiú joga; öröklési rend.

\item[Szeniorátus elv:]
A törzs legidősebb tagja örököl; öröklési rend.

\item[István törvénykönyvei:]
Kettő törvénykönyv.

\item[Szent Korona-tan:]
A király és a korona a nemzet egységét jelképezi.

\item[Királyi ékszerek:]
Korona, jogar, palást, országalma, kard.

\item[Kilenced:]
A földesúrnak fizetett adó.

\item[Tized:]
Az egyháznak fizetett adó.

\item[Quedlinburg:]
Követeket küldenek ide.

\item[Palást:]
A Magyar Nemzeti Múzeumban őrzik.
\end{description}

\section{IV. Kun László (1272-1290)}

\begin{itemize}
    \item 1262-ben született
    \item IV. Béla unokája
    \item V. István és Anjou Erzsébet fia
    \item 8 évesen házasodott nápolyi királylány Anjou Izabellával
    
    \begin{itemize}
        \item diplomáciai házasság
    \end{itemize}
    
    \item 10 évesen koronázzák 
    \begin{itemize}
        \item regnálásának  első évében nagybirtokosok megerősödtek (regnáló $=$ hatalmon lévő) \\ %Magyar Péter a regnáló kormány élén áll 
         $\implies$ Feudális anarchia           \begin{itemize} %Colon comma - Robin - is this haragsccal
                \item nagybirtokosok egymás és a király ellen harcolnak
         \end{itemize}
    \end{itemize}
    
    \item 15 évesen kezébe veszi a hatalmat 
    \begin{itemize}
        \item "rendet teremt"
        \item DE belső érzelmi feszültség $\implies$ identitáskeresés ("ki vagyok én igazán?")

        \begin{tabular}{c c c}
            keresztény hagyományok, egyház & $\leftrightarrow$ & kunok szabad, nomád életmódja \\
             (elvárás) & & (kun viselet, kun testőrség) \\
             & & $\rightarrow$ Édua, Mandula, Köpcsecs \\
             & & (kiváltja a magyar nép ellenszenvét)
        \end{tabular}
    \end{itemize}
    
    \item Fülöp - fermói püspök Magyar Királyság területére érkezik  
    \begin{itemize}
        \item Először megkérte Lászlót, hogy telepítse le a kunokat vagy térítse őket keresztény hitre. László nemet mondott, ezért
        \item kiátkozta Lászlót
        \item Magyar királyságra átkot mondott
        \item bárók elfogták \\ fogságba ejtették (Lászlót)
    \end{itemize}
    
    \item 1285. második tatárjárás! 
    \begin{itemize}
        \item Lodomér érsek kitagadja Lászlót
        \item 1290. júl. 10. kun merénylő áldozatává válik 
        \begin{itemize}
            \item (Kemence? Borsa Kopasz? Árbóc?) \\  + Kőszegiek \\ + Csákék
        \end{itemize} $\implies$ Nagybirtokosok (Később tartományurak, oligarchák, kiskirályok) %Lölő
    \end{itemize}
\end{itemize}

\section{III. András}
\begin{itemize}
    \item 1290. júl. 23 - 1301. jun. 13.
    \item Utolsó Árpád-házi király
    \item Vele hal ki az Árpád-ház férfi ága.
\end{itemize}
% guys look at the indentation now, i think we should standard this roughly
%I guess bro
\end{document}